\vssub
\subsection{~程序文件} \label{sec:files}
\vssub

\ws\ 源代码文件存放在扩展名为 {\file ftn} 的文件中\footnote{~少数由他人
提供的模块除外。}。从 2.00 版开始，代码按模块组织。只有主程序没有封装
到模块中。最初，变量与代码模块捆绑在一起，形成单一的静态数据结构。在
模式 3.06 版中，引入了独立的动态数据结构，允许单个程序中存在多个波浪
网格，为开发多网格模式驱动器做准备。

下文会较详细地说明各模块包含的子程序。各种子程序之间的关系以图形方式
表示在图~\ref{fig:w3init} 和图~\ref{fig:w3wave} 中。这里考虑三组代码。
第一组是主波浪模式子程序模块，通常可由文件名结构
{\file w3{\it{xxxx}}md.ftn} 识别。这些模块见 \para\ref{sec:wave_mod}。
第二组是多网格波浪模式驱动器专用模块，通常可由文件名结构
{\file wm{\it{xxxx}}md.ftn} 识别。这些模块见 \para\ref{sec:multi_mod}。
最后一组由辅助程序和波浪模式驱动器组成，见 \para\ref{sec:aux_mod}。
第~\ref{sec:data_ass} 节简要说明数据同化模块。
